\documentclass[11pt,a4paper]{letter}
\usepackage[utf8]{inputenc}
\usepackage[T1]{fontenc}
\usepackage{mlmodern}
\usepackage{amsmath}
\usepackage{amsfonts}
\usepackage{amssymb}
\usepackage{graphicx}
\usepackage{geometry}
\usepackage{csquotes}
\usepackage[final]{microtype}
\geometry{margin=1in}

\address{Aniara Lyn Lauron \\ Gåseholmvej 61 \\ Herlev, Denmark \\ Email: aniaralyn.lauron@gmail.com \\ Phone: +45 91 80 63 94}
\date{February 22, 2026}

\begin{document}

\begin{letter}{Admissions Committee \\ VIA University College \\ Banegårdsgade 2, 8700 Horsens \\ Denmark}

\opening{Dear Admissions Committee,}

I am writing to express my strong interest in the Bachelor of Engineering in Software Technology programme at VIA University College, starting in the fall of 2026. I hold a Bachelor of Science in Pharmacy from the University of Southern Philippines and am eager to transition into software development, where I can combine my scientific foundation with creative engineering, sustainable thinking, and modern digital tools. The Software Technology Engineering programme promises exactly the kind of hands on, project based education I am looking for starting with programming, web development and software design, and later exploring algorithms, DevOps, machine learning and embedded systems. I am motivated to build software from scratch and to learn through doing, and I appreciate that the programme welcomes students without prior coding experience but with a strong logical and mathematical mindset.

My academic foundation in pharmacy has equipped me with a solid understanding of the scientific method, data analysis, and ethical responsibility. During my studies I conducted research on the anti-inflammatory effects of natural extracts, culminating in my thesis \enquote{The Effect of Hagonoy Leaves Extract on the Inflamed Paws of Male Albino Rabbits}. The project involved chemical extraction, animal experimentation, and statistical evaluation, experiences that honed my analytical thinking and gave me an appreciation for modelling complex systems. Over time I became fascinated by how software models, algorithms and systems can similarly describe and solve real world problems.

Beyond academics, I have gained practical experience in diverse roles. As a pharmacy technician at Optum Global Solutions Inc. and Collabera Technologies Private Limited I managed medication dispensing and collaborated with healthcare teams, tasks that required precision, reliability, and teamwork. My freelance work as an English teacher helped develop my communication skills, and my current role as a hotel assistant at The Ellen Group has enhanced my problem solving and customer service abilities. Each position taught me how to approach problems methodically and adapt to new tools and workflows, qualities that are essential for a software engineer working on collaborative projects.

Being surrounded by engineers and technologists during my international experiences sparked my passion for technology. Living in the Philippines, South Korea, the Netherlands, and Denmark has shaped my adaptability and global perspective. The Philippines fostered resilience and community values, Korea's dynamic tech culture inspired innovation, the Netherlands' focus on sustainability motivated my desire to apply computation to real challenges, and Denmark's collaborative, pedagogy driven approach aligns with how I learn best. I am driven by a genuine desire to help others and to continuously improve, qualities that will serve me well in the collaborative, project based environment at VIA Horsens.

To begin bridging my background with software, I have pursued hands on technical projects. In 2024 I designed and built an Arduino based automated plant irrigation system, programming firmware in C++ to automate watering schedules and integrating sensors and actuators. The project taught me the basics of input/output, control loops and embedded programming, and it was rewarding to see working software control a physical system. The IoT lab and embedded tracks at VIA are exactly the kinds of environments I wish to explore further.

Additionally, my 2026 custom piano computer workstation project involves woodworking, design, and assembly, reflecting creativity and precision in crafting functional solutions. My 2023 silver ring making project demonstrated metalworking and engraving skills, highlighting my flexibility and capacity to learn new disciplines. These diverse hobbies reinforce my belief that engineering and programming are crafts built through experimentation and iteration.

The Software Technology Engineering programme appeals to me for its structured progression from single user to client/server to heterogeneous systems, culminating in large-team collaborative development and industry internship. I value the emphasis on sustainability, digitalisation, innovation and entrepreneurship that runs through every semester, and I look forward to the opportunity to specialise in areas such as interactive media, business information systems, or machine learning and AI engineering. The combination of classroom theory, lab work in IoT and XR, company visits, guest lecturers and an internship will prepare me to design, develop and operate software that makes a difference.

Looking ahead, after completing the B.Eng. in Software Technology, I plan towork at the intersection of software and life sciences. My long term goal is to develop tools and systems, perhaps for companies like Novo Nordisk or Coloplast, that improve healthcare and quality of life through intelligent software. VIA's focus on innovation and company collaboration makes it the perfect springboard toward these ambitions.

I am excited about the opportunity to join VIA Horsens and contribute my diverse perspective to the Software Technology Engineering programme. Thank you for considering my application. I look forward to the possibility of discussing how my background and aspirations align with VIA's mission.


\closing{Sincerely, \\
Aniara Lyn Lauron}

\end{letter}

\end{document}